\addtocontents{toc}{\protect\Needspace{6\baselineskip}}
\chapter{服务功能}\label{ch:svc}

RTC、SNMP 和 615A 均以 UDP 报文工作，经服务通道 \code{afdx_svc} 与应用共享收发核心。服务报文按配置的 VL 参与 BAG 调度、完整性检查和冗余管理。


\section{服务通道}\label{sec:svcch}

\code{afdx_svc} 由组合选择和五个状态寄存器组成：发送侧为 \code{busy}、\code{owner}、\code{turn}，接收侧为 \code{rx_active}、\code{rx_claimed}，结构见图~\ref{fig:svc}。

\figN[服务通道]{f26_svc}{服务通道：(a) 应用与服务的整消息发送仲裁；(b) 按 UDP 目的端口进行接收分发}{fig:svc}

应用口与三个服务发送口按整条消息轮询，次序为应用、RTC、SNMP、615A。选中的来源独占消息通路，直至末字节握手，随后从下一来源开始选择。服务发送口均空闲时，应用的数据和 ready 组合直通发送核心。发送仲裁器由三个状态寄存器构成：\code{busy} 指示当前消息传输忙，\code{owner} 锁存当前来源，\code{turn} 指示下一轮轮询起点；末字节握手时 \code{busy} 清零、\code{turn} 指向当前来源的下一个。由此任一来源在提出请求后，至多等其余三个来源各传完一条消息。

接收核心的交付口同时广播给应用和三个服务。各服务按 UDP 目的端口给出认领信号，首字节的认领结果保持到末字节：被认领的消息只送服务，应用口的 valid 和 last 为 0；其余消息原样送给应用。三个服务使用不同的 UDP 端口，且都不对交付口施加反压。认领信号由各服务按 \code{d_dport} 组合产生，首字节当拍直接决定应用口的 \code{valid}，之后由寄存的 \code{rx_claimed} 保持。\code{rx_active} 选择当前认领或寄存的认领结果，其反相信号与 \code{d_valid} 相与形成应用口有效信号。

\section{RTC 时钟同步}\label{sec:rtc}

\code{afdx_rtc} 维护 32 位秒加 32 位纳秒的实时时钟，每拍前进 8~ns，纳秒满 $10^9$ 向秒进位；秒的起点为 2000-01-01 00:00:00 UTC，线上 NTP 秒为其加 3155673600。同步采用 RFC 4330 的 SNTPv4 主从方式，交互见图~\ref{fig:rtc}。

\figN{f14_rtc}{SNTP 主从对时交互及四个时间戳的采样位置}{fig:rtc}

\subsection{时间戳采样与收发时序}

\code{afdx_rtc} 包含本地时钟、一个 48 字节接收移位寄存器、发送字节状态机、校时计算状态机和三个 \code{afdx_ntp_fraction} 换算器（发送时间戳编码、接收或参考时间戳编码、回应时间戳解码），两个状态机见图~\ref{fig:rtcfsm}。

从端每隔 \code{PERIOD} 发一个请求，记下发送时刻 $t_1$；主端在收到请求时记 $t_2$，回应发送时记 $t_3$，并把请求的 Transmit 时间戳逐位复制到回应的 Originate 字段；从端收到回应时记 $t_4$。从端只接受 Originate 与本次请求逐位相等的回应，随后计算
\begin{equation}
\text{offset} = \frac{(t_2-t_1)+(t_3-t_4)}{2}
\end{equation}
并加到当前时钟。一个周期内未收到有效回应时，下一周期以新的 Transmit 时间戳重发，直至校时成功。

发送时间戳取消息首字节与发送核心的握手时刻，接收时间戳取核心交付首字节的时刻。时间戳在握手后开始换算，并在对应字段发出前完成。

\figN[RTC 的两个状态机]{f27_rtc}{RTC 状态机：(a) 报文发送；(b) 回应时间戳解码与时钟校正}{fig:rtcfsm}

主端有待回应请求或从端倒计时归零时启动发送。无反压时每字节用 4 拍：TX\_GROUPS 先在两组字段中各选一个 32 位字，TX\_WORD 选字，TX\_BYTE 选字节并置 \code{tx_valid}，TX\_HANDSHAKE 等 \code{tx_ready}。一个 48 字节报文至少 192 拍。Transmit 字段位于第 40--47 字节，距首字节握手至少 160 拍，换算耗时 65 拍。主端的 Receive 字段 $t_2$ 在请求首字节到达时开始换算，回应在整条请求收完后开始发送。

\subsection{时间换算与偏差校正}

\code{afdx_ntp_fraction} 负责 NTP 的 32 位二进制小数与本地纳秒值之间的换算。纳秒编码为 NTP 小数时，用 32 轮、每轮 2 拍的逐位比较减法计算 $\lfloor \text{ns}\times 2^{32}/10^9 \rfloor$；反方向用逐位加与右移并在末轮四舍五入。每拍最多一级 32 位加减，整纳秒时间戳往返换算得到原值。

秒差和纳秒差分开处理，各拍至多一级 32 位运算，步骤见图~\ref{fig:rtcfsm}(b)：先依次解码回应中的 $t_2$、$t_3$ 小数部分（各 65 拍），再用 9 拍计算出两段差、求和、除以 2，并把奇数秒的半秒折入纳秒、处理借位进位。校正量在 ADD\_OFFSET 步加至当前时钟，至 APPLY 步写回共经过 4 拍，故在 LATENCY 步补偿 32~ns；一次校时约 140 拍，期间时钟正常推进。

本端正在等待回应且发送状态机空闲时，以版本 VN=4、模式 mode=4 及 Originate 与本次请求 Transmit 逐位相等作为校时计算的启动条件，服务器 stratum 与全零时间戳不参与判定。迟到的旧回应因 Originate 不符而被丢弃。计算期间或 \code{rx_done} 有效的周期不发起新请求。

\subsection{报文字段与寄存器}

请求与回应均为 48 字节，所有多字节字段高字节先发，见表~\ref{tab:sntp}。双方 UDP 源、目的端口均为 \code{UDP_PORT}（默认 123）。

\begin{table}\centering\small
\caption{SNTPv4 报文字段}\label{tab:sntp}
\begin{tabularx}{\linewidth}{>{\centering\arraybackslash}p{1.4cm}LL}
\toprule
字节 & 请求（从端） & 回应（主端） \\
\midrule
0 & \code{0x23}：LI=0、VN=4、mode=3 & \code{0x24}：LI=0、VN=4、mode=4 \\
1 & stratum=0 & stratum=1 \\
2--3 & Poll=4、Precision=$-26$ & 同左（约 14.9~ns） \\
4--11 & Root Delay、Root Dispersion 均为 0 & 同左 \\
12--15 & 0 & Reference ID=\code{LOCL} \\
16--23 & 0 & 最近一次配置装入的时刻 \\
24--31 & Originate=0 & 请求 Transmit 的逐位复制 \\
32--39 & Receive=0 & $t_2$ \\
40--47 & Transmit=$t_1$ & $t_3$ \\
\bottomrule
\end{tabularx}
\end{table}

RTC 寄存器位于 \code{0x80}--\code{0x88}，见附录表~\ref{tab:regs}。写时间先写 \code{SEC}、\code{NS} 两个暂存字，再以 \code{CTRL} 的装入位一次生效；读 \code{SEC} 时同时锁存纳秒值，随后读 \code{NS} 得到同一时刻的值。偏差以 \code{OFFSET_SEC}$\times 10^9 +$\code{OFFSET_NS} 纳秒表示。


\section{SNMP 网络管理}\label{sec:snmp}

\code{afdx_snmp} 是服务通道上的 SNMP 代理，认领 UDP 目的端口 161，回应的目的端口为请求的源端口，经配置的发送 VL 送出。图~\ref{fig:snmp} 是其处理流程。

\figN{f15_snmp}{SNMP 代理处理流程}{fig:snmp}

\subsection{协议对象与操作}

MIB 包含 35 个只读对象与 1 个读写对象（\code{USER_ID}），共 36 个。对象分为 sysDescr、sysUpTime 两个系统对象，以及演示用私有子树 \code{1.3.6.1.4.1.424242.3} 下的 34 个对象，覆盖版本号、\code{USER_ID}、RTC 秒和同步次数、全部 27 个收发统计计数器、SNMP 请求数、community 错误数和 Trap 发送数。对象值直接取自寄存器。

\Needspace{4\baselineskip}
代理支持以下协议版本和操作：
\begin{itemize}
  \item \textbf{版本：}SNMPv1（version=0）和 v2c（version=1）。
  \item \textbf{Get / GetNext：}Get 精确匹配标量 OID；GetNext 按子标识符逐个比较，返回严格大于请求 OID 的第一个对象，可从 \code{1.3} 遍历全部对象。对象或实例未命中时，v1 回应 noSuchName，v2c Get 统一回应 noSuchObject；v2c GetNext 遍历至末尾时回应 endOfMibView。
  \item \textbf{Set：}\code{USER_ID} 为读写配置对象，类型为 Gauge32，取值 0--65535。代理先检查 OID、类型和取值范围，再经配置总线写入并读回，以读回值组装回应。只读对象、类型不符和越界分别按 v1/v2c 的错误码回应，配置不变。
  \item \textbf{Trap：}发送 SNMPv2-Trap（PDU \code{0xA7}）到 UDP 端口 162，含 sysUpTime.0 和 snmpTrapOID.0 两个变量绑定。\code{TRAP_ENABLE} 已置位且 \code{ENABLE} 从 0 变为 1 时发 coldStart；community 不符且 \code{TRAP_ENABLE} 已置位时发 authenticationFailure。
\end{itemize}
每个请求按单变量绑定处理；community 不符时不回应并计数 \code{BAD_COMMUNITY}。

\subsection{请求解析与回应组装}

代理的主状态机有 44 个状态，按功能分为 7 组，见图~\ref{fig:snmpfsm}。请求先完整写入 256$\times$8 的请求 RAM，解析与查找通过公共读取状态序列访问 RAM：原状态分支将返回位置写入 \code{return_state}，读取序列用 4 拍完成越界检查、发出读请求、等待同步读出和锁存数据，随后转回原状态分支。\code{Fetch} 状态组由解析、\code{Lookup}、\code{Value} 和 \code{Encode} 共用；每个解析状态只处理一个已就绪的字节，单拍逻辑只有一级比较。Set 未命中时经 Value 组的 \code{SET_ERROR} 编码错误回应，异常处理完成后返回接收状态。

\figN[SNMP 代理状态机（按功能分组）]{f28_snmpfsm}{SNMP 代理的请求解析、对象访问与回应组装状态机}{fig:snmpfsm}

代理按 BER 格式依次解析 SEQUENCE、version、community 和 PDU；PDU 内依次是 request-id、error-status、error-index 和变量绑定表。长度字段支持短形式和 \code{0x81}、\code{0x82} 两种定长形式。回应在 256$\times$8 的回应 RAM 中逐字节组装，原样保留 request-id 和 community，再经服务发送口输出。

对象查找按对象表顺序逐个尝试 36 个候选，每个候选按子标识符逐段比较：多字节子标识符按 BER 的 7 位分组累加还原成整数后比较数值，使 GetNext 按 OID 的数值顺序选择对象。

LEN1--LEN4 四拍自内向外计算出值、变量绑定、绑定表、PDU 和整条消息的长度，BUILD 再按“类型、长度、内容”逐字段写入回应 RAM；community、OID 和回显的值用 COPY 状态从请求 RAM 逐字节复制。回应写完后由发送状态逐字节读出，每字节 4 拍送入服务发送口。

代理一次处理一个请求，只在 IDLE 接收新请求；处理期间到达的请求仍被认领并计入请求数，但不回应，由管理端按超时重发。超过 256 字节的请求整条丢弃。

\subsection{配置访问与寄存器}

代理经与主机共享的配置读通路访问寄存器：主机访问其他地址时主机优先，同一地址的主机读和代理读在同一拍合并；代理写只在主机无读写且 \code{cfg_ready} 为 1 时授予。

SNMP 寄存器位于 \code{0x90}--\code{0xA3}，包括启用、回应 VL、community 长度和内容、请求计数、Trap 开关与计数及 sysDescr，见附录表~\ref{tab:regs}。

\section{615A 数据加卸载}\label{sec:a615}
\begingroup
\setlength{\parskip}{0pt}

\code{afdx_615a} 实现 ARINC 615A 的目标端：加载端经 TFTP 送来上传清单，目标主动读取头文件和数据文件，校验后把数据留在 8~KiB 加载区，并通过状态文件报告进度与结果；目标也能以介质定义卸载把加载区作为文件 \code{MEMORY.BIN} 读出。协议文件版本默认 \code{A4}，兼容 \code{A3}；目标标识为 \code{DEMC3_1}。

设计范围为单文件上传和介质定义卸载，信息操作与操作员定义卸载在初始化阶段拒绝；FIND、WAIT、ABORT、整文件重试和多文件加载不在实现范围内。TFTP 采用固定 512 字节块长，请求中的 RFC 2347 选项按忽略处理。ERROR 报文按目的 TID 关联会话，来源端口不参与校验。请求按操作码和零结尾字段解析，名称与模式字段的读取不以实际接收长度为终止条件。

\subsection{报文分派与异常处理}

\code{afdx_615a} 的主状态机约 70 个状态，按图~\ref{fig:a615fsm} 分组。收到的报文先整条存入 1~KiB 的接收缓冲，由分派状态 DISPATCH 按操作码、端口（TID）和块号分派给六类处理过程。图中的 Transfer/status 包含 List、Header、Data、Download 和 Status 处理；End 完成后返回等待状态。

各处理过程共用字节读取、整数拼接与报文输出状态序列。读字节以 \code{return_state} 保存返回位置，读整数另用 \code{num_return} 保存字段处理的返回位置；报文输出在发送 RAM 中逐字节组包，发送后按返回状态寄存器 \code{after_tx} 转移，通常返回等待报文状态 \code{WAIT_PACKET}。这些公共序列使各处理分支共用相同的同步读与发送时序。

数据或文件检查失败时，锁存错误码 \code{reason} 并转入 \code{FAILURE} 状态，依次发送 TFTP ERROR 和失败状态文件；状态文件传输失败（含重试耗尽）时结束。请求名称或模式错误只回 ERROR；LCI/LNO 在初始化阶段拒绝。

\figN[615A 目标状态机（按功能分组）]{f29_615a}{615A 目标端的请求分派、文件传输与状态报告状态机}{fig:a615fsm}

接收缓冲一次保存一个报文。当前报文未取走时，新到报文仍被认领，但不写入缓冲；对端未收到回应时，按超时机制重发。

\subsection{TFTP 传输与重传}

TFTP 使用 RRQ、WRQ、DATA、ACK、ERROR 五种报文，octet 模式，块长 512 字节，块号从 1 开始；WRQ 以 ACK0 建立传输。请求目的端口为 \code{REQUEST_PORT}（默认 59），建立传输后双方使用各自的传输标识符（TID），目标 TID 从 \code{0x8100} 递增。主文件与状态文件分别保存 TID 和块号。

报文在 RAM 中保留到回应到达；超时（默认 2~s）后原样重发上一报文，重传上限默认 1 次，耗尽以原因 5 失败。重复 DATA 只重发 ACK，不重复写 RAM。

\subsection{上传与文件校验}

上传流程见图~\ref{fig:a615}(a)：
\begin{enumerate}
  \item 加载端读取 \code{DEMC3_1.LUI}，目标返回接受状态 0001。
  \item 目标写 \code{DEMC3_1.LUS} 报告初始状态。
  \item 加载端写 \code{DEMC3_1.LUR} 上传清单，含头文件名和件号；版本不支持则以原因 1 拒绝。
  \item 目标发 LUS 0002 后读取清单中的 LUH 头文件；头文件校验通过后读取其列出的数据文件。
  \item 数据块按序写入加载区并累计 CRC。首个满块写入后暂缓 ACK，先经独立传输发送一次进度 LUS，再确认数据块。
  \item 末块（短块或空块）到达后，长度和 CRC 均正确才确认，并发送 LUS 0003、进度 100；否则发送 TFTP ERROR 和 LUS 1003。
\end{enumerate}

\figN[615A 交互]{f16_615a}{615A 单文件上传与介质定义卸载交互}{fig:a615}

两种传输的末块均少于 512 字节；文件长度为 512 字节的整倍数时，补发零长度 DATA 块。上传在 LUH 校验完成后先报告进度，再请求数据文件。

LUH 支持 ARINC 665-1 的 8002 版本：读取头文件字数、数据文件列表指针、数据文件数、文件名、数据文件字数和 CRC-16，检查头文件实际长度、版本、列表数量、头文件 CRC 及数据文件的长度和 CRC。件号和硬件 ID 不参与匹配；支持文件、用户数据、扩展点、Load CRC-32、解压和解密不参与当前加载处理。

CRC-16 的多项式为 \code{0x1021}，最高位优先、初值 \code{0xFFFF}、结果不取反，以 8 位并行的组合函数每拍处理一字节。

新 LUS 计数加 1，重传不加。进度取已写字节占预期总字节的整数百分比，成功前上限为 99。计算时先由移位相加完成乘 100，再以预期总字节数为减数逐拍迭代求商，最后通过减 10 迭代分离十位与个位，转换为十进制 ASCII 字符。

\subsection{介质定义卸载}

卸载流程见图~\ref{fig:a615}(b)。加载端读 \code{DEMC3_1.LND}，目标回应后发初始 LNS；加载端写 LNR 列表，只接受 \code{MEMORY.BIN}；目标发 LNS 0002 后主动写该文件，按 512 字节块停等发送最近一次成功上传的内容，首个满块确认后报告进度，末块确认后发 LNS 0003。卸载只读加载区，不改变已加载字节数。

上传和卸载互斥，操作在最终状态文件获 ACK 后才记为成功；新上传开始时撤销旧的有效长度。加载区采用 8~KiB 片上 RAM，供主机经配置总线按字节读回或用于介质定义卸载。


\subsection{存储与寄存器}

五块 RAM 分别存放接收报文（1~KiB）、发送与重传报文（1~KiB）、LUR（512~B）、LUH（2~KiB）和加载数据（8~KiB），均以 \code{afdx_ram_sdp} 例化。

寄存器 \code{0xB0}--\code{0xBB} 见表~\ref{tab:regs}。\code{STATUS} 报告空闲、进行中、成功、失败及原因码：0 无错误；1 元数据不支持或无效（版本、文件名、数量）；2 容量或长度超限；3 CRC 错；4 对端 TFTP ERROR 或模式错误；5 重传耗尽。加载区由 \code{MEMORY_ADDRESS} 选字节、\code{MEMORY_DATA} 读回。
\endgroup
